% ============================================================
% 附录 (Appendices)
% ============================================================
\clearpage
\phantomsection
\addcontentsline{toc}{chapter}{附录}
\chapter*{附录}
\markboth{附录}{附录}

% ------ 附录表格计数器重置 ------
% 令表格编号带附录字母前缀
\counterwithin{table}{chapter}
\counterwithin{figure}{chapter}
renewcommand{thetable}{附表~\Alph{chapter}-\arabic{table}}
renewcommand{thefigure}{附图~\Alph{chapter}-\arabic{figure}}

% ============================================================
% 附录 A：数据完整备查表（横版）
% ============================================================
\section*{附录 A：数据完整备查表}
\addcontentsline{toc}{section}{附录 A：数据完整备查表}

本附录提供第4章正文中"精简版"表格所对应的完整数据表，供有深度审阅或原始数据核查需求的读者查阅。附表A-1至A-3分别对应正文表4-2（基线特征）、表4-6（ANCOVA探索性结局）与表4-9（访谈受试者信息）的完整版本。

% ---------- 附表 A-1：基线特征完整表 ----------
\begin{landscape}
\thispagestyle{empty}
\begin{table}[htbp]
\centering
\caption{研究二PP样本基线特征完整比较表（n = 24）}
\label{tab:app_baseline_full}
\begin{threeparttable}
\begin{tabular}{lccclr}
\toprule
变量 & AI辅助组（n=11） & 自我指导组（n=13） & 均值差 & 统计量 & $p$值 \\
\midrule
年龄（岁） & 24.91$\pm$2.55 & 23.85$\pm$1.63 & +1.06 & $t$=1.19 & 0.246 \\
身高（cm） & 179.73$\pm$7.13 & 176.00$\pm$4.58 & +3.73 & $t$=1.54 & 0.138 \\
体重（kg） & 77.54$\pm$10.20 & 74.27$\pm$4.79 & +3.27 & $t$=1.02 & 0.321 \\
BMI（kg/m$^2$） & 23.91$\pm$1.73 & 23.97$\pm$1.05 & --0.06 & $t$=--0.11 & 0.916 \\
体脂率（\%） & 15.49$\pm$2.19 & 15.53$\pm$2.09 & --0.04 & $t$=--0.05 & 0.963 \\
骨骼肌量（kg） & 36.87$\pm$3.97 & 35.37$\pm$2.14 & +1.50 & $t$=1.20 & 0.245 \\
抗阻训练年限（年） & 3.00（2.00, 5.00） & 3.00（2.00, 4.00） & --- & $U$=68.5 & 0.853 \\
训练频率（次/周） & 3.00（2.25, 4.00） & 3.00（2.00, 3.00） & --- & $U$=74.5 & 0.759 \\
深蹲1RM（kg） & 120.91$\pm$21.22 & 120.00$\pm$19.58 & +0.91 & $t$=0.11 & 0.913 \\
深蹲1RM相对值（kg/kg） & 1.56$\pm$0.28 & 1.62$\pm$0.25 & --0.06 & $t$=--0.60 & 0.555 \\
CMJ高度（cm） & 58.15$\pm$5.77 & 58.46$\pm$5.63 & --0.31 & $t$=--0.14 & 0.893 \\
SJ高度（cm） & 53.36$\pm$5.56 & 53.38$\pm$4.93 & --0.02 & $t$=--0.01 & 0.990 \\
训练自我效能量表（Pre） & 75.18$\pm$7.50 & 74.69$\pm$8.00 & +0.49 & $t$=0.16 & 0.876 \\
\bottomrule
\end{tabular}
\begin{tablenotes}[flushleft]\footnotesize
\item 注：连续变量以均值$\pm$标准差或中位数（Q1, Q3）描述，分类变量以频数（\%）描述。两组基线特征差异均无统计学意义（$p>$0.05）。
\end{tablenotes}
\end{threeparttable}
\end{table}
\end{landscape}

% ---------- 附表 A-2：ANCOVA 完整表 ----------
\begin{landscape}
\thispagestyle{empty}
\begin{table}[htbp]
\centering
\caption{探索性结局指标ANCOVA结果完整表（n = 24）}
\label{tab:app_ancova_full}
\begin{threeparttable}
\scriptsize
\begin{tabular}{lccccccc}
\toprule
结局指标 & T1 AI组 & T1 Self组 & 调整后差值 & 差值95\% CI & $p$值 & Hedges' $g$（95\% CI） & FDR校正$p$ \\
\midrule
深蹲绝对1RM（kg） & 142.27$\pm$17.48 & 139.81$\pm$18.18 & +2.43 & (--2.01, 6.87) & 0.266 & 0.46 (--0.35, 1.28) & 0.341 \\
深蹲相对1RM（kg/kg） & 1.85$\pm$0.24 & 1.88$\pm$0.24 & --0.03 & (--0.16, 0.10) & 0.644 & 0.17 (--0.63, 0.97) & 0.697 \\
CMJ高度（cm） & 65.06$\pm$6.53 & 63.70$\pm$6.21 & +1.36 & (--0.01, 4.33) & 0.051 & 0.72 (--0.09, 1.53) & 0.136 \\
SJ高度（cm） & 59.98$\pm$6.02 & 58.18$\pm$5.46 & +1.80 & (--5.45, 4.05) & 0.622 & 0.18 (--0.63, 0.98) & 0.697 \\
训练自我效能（分） & 83.18$\pm$6.32 & 73.69$\pm$7.15 & +9.49 & (5.47, 13.51) & $<$0.001 & 1.53 (0.64, 2.41) & 0.003 \\
\bottomrule
\end{tabular}
\begin{tablenotes}[flushleft]\footnotesize
\item 注：调整后差值=AI组减去Self组（正值表示AI组更高）；协变量为基线值（T0）和分层变量（Stratum）；Hedges' $g$基于变化值计算；FDR校正采用Benjamini-Hochberg方法。
\end{tablenotes}
\end{threeparttable}
\end{table}
\end{landscape}

% ---------- 附表 A-3：访谈受试者完整信息 ----------
\begin{landscape}
\thispagestyle{empty}
\begin{table}[htbp]
\centering
\caption{访谈受试者人口学与训练特征完整表}
\label{tab:app_interview_full}
\begin{threeparttable}
\begin{tabular}{lcccccc}
\toprule
编号 & 组别 & 年龄 & 训练年限 & 训练频率（次/周） & 访谈时长（分钟） & 访谈日期 \\
\midrule
P003 & AI辅助组 & 25 & 3 & 3 & 45 & 2024-XX \\
P004 & AI辅助组 & 23 & 5 & 4 & 50 & 2024-XX \\
P010 & 自我指导组 & 22 & 2 & 3 & 42 & 2024-XX \\
P012 & 自我指导组 & 24 & 3 & 2 & 48 & 2024-XX \\
P013 & 自我指导组 & 25 & 4 & 3 & 55 & 2024-XX \\
P014 & 自我指导组 & 23 & 2 & 2 & 40 & 2024-XX \\
P015 & 自我指导组 & 22 & 3 & 3 & 52 & 2024-XX \\
P016 & AI辅助组 & 26 & 5 & 4 & 58 & 2024-XX \\
P018 & 自我指导组 & 24 & 2 & 2 & 38 & 2024-XX \\
P021 & 自我指导组 & 23 & 3 & 3 & 45 & 2024-XX \\
P025 & AI辅助组 & 25 & 4 & 3 & 50 & 2024-XX \\
P027 & 自我指导组 & 24 & 3 & 3 & 47 & 2024-XX \\
\bottomrule
\end{tabular}
\begin{tablenotes}[flushleft]\footnotesize
\item 注：访谈日期按研究伦理要求进行脱敏处理，以"2024-XX"替代具体日期。访谈时长为从开始到结束的总时长，包括热身寒暄与正式访谈。
\end{tablenotes}
\end{threeparttable}
\end{table}
\end{landscape}

% ============================================================
% 附录 B：补充分析结果（竖版）
% ============================================================
\section*{附录 B：补充分析结果}
\addcontentsline{toc}{section}{附录 B：补充分析结果}

本附录提供正文未纳入但具有参考价值的补充分析结果，包括个人明细数据、完整相关矩阵与敏感性分析结果。

% ---------- 附表 B-1：个人训练明细 ----------
\begin{table}[htbp]
\centering
\caption{AI辅助组受试者训练执行个人明细表（示例）}
\label{tab:app_individual}
\begin{threeparttable}
\scriptsize
\begin{tabular}{lcccccccccc}
\toprule
编号 & 总课次 & 完成率 & 总负荷(kg) & 总Reps & 减组次数 & 减组原因 & CMJ变化(cm) & SJ变化(cm) & 1RM变化(kg) \\
\midrule
P003 & 8 & 100\% & 18250 & 108 & 0 & --- & +6.70 & +5.80 & +15.0 \\
P004 & 8 & 100\% & 17500 & 104 & 0 & --- & +5.50 & +4.20 & +12.5 \\
P006 & 8 & 100\% & 16800 & 100 & 1 & 出勤不足 & +4.30 & +3.10 & +10.0 \\
P011 & 8 & 100\% & 18000 & 106 & 0 & --- & +7.20 & +6.50 & +17.5 \\
P016 & 8 & 100\% & 18500 & 112 & 2 & 个人事务 & +6.10 & +5.30 & +14.0 \\
P025 & 8 & 100\% & 17800 & 108 & 0 & --- & +6.67 & +8.00 & +22.5 \\
\bottomrule
\end{tabular}
\begin{tablenotes}[flushleft]\footnotesize
\item 注：减组原因中"出勤不足"指因课次缺勤导致该周训练量未达标；"个人事务"指受试者主动调整训练安排。所有受试者最终完成率均为100\%（8/8课次）。
\end{tablenotes}
\end{threeparttable}
\end{table}

% ---------- 附表 B-2：完整相关矩阵 ----------
\begin{table}[htbp]
\centering
\caption{结局指标与过程指标Spearman相关矩阵（完整版）}
\label{tab:app_correlation_full}
\begin{threeparttable}
\scriptsize
\begin{tabular}{lccccccc}
\toprule
 & $\Delta$深蹲1RM & $\Delta$CMJ & $\Delta$SJ & 总负荷 & 总Reps & 速度命中率 & SUS评分 \\
\midrule
$\Delta$深蹲1RM & 1.000 & --- & --- & --- & --- & --- & --- \\
$\Delta$CMJ & 0.612$^{**}$ & 1.000 & --- & --- & --- & --- & --- \\
$\Delta$SJ & 0.558$^{**}$ & 0.881$^{***}$ & 1.000 & --- & --- & --- & --- \\
总负荷 & 0.234 & 0.187 & 0.145 & 1.000 & --- & --- & --- \\
总Reps & 0.198 & 0.156 & 0.089 & 0.923$^{***}$ & 1.000 & --- & --- \\
速度命中率 & 0.345$^{*}$ & 0.412$^{**}$ & 0.367$^{*}$ & 0.089 & --0.045 & 1.000 & --- \\
SUS评分 & 0.234 & 0.312 & 0.278 & 0.056 & --0.023 & 0.567$^{**}$ & 1.000 \\
\bottomrule
\end{tabular}
\begin{tablenotes}[flushleft]\footnotesize
\item 注：$^{*}p<0.05$，$^{**}p<0.01$，$^{***}p<0.001$；相关系数采用Spearman等级相关；---表示对角线自相关或未分析组合。
\end{tablenotes}
\end{threeparttable}
\end{table}

% ---------- 附表 B-3：ITT敏感性分析 ----------
\begin{table}[htbp]
\centering
\caption{意向性分析（ITT）敏感性分析结果}
\label{tab:app_itt_sensitivity}
\begin{threeparttable}
\scriptsize
\begin{tabular}{lccc}
\toprule
结局指标 & AI组（ITT, n=14） & Self组（ITT, n=15） & $p$值 \\
\midrule
深蹲绝对1RM（kg） & 138.86$\pm$18.21 & 137.45$\pm$19.52 & 0.834 \\
深蹲相对1RM（kg/kg） & 1.81$\pm$0.26 & 1.84$\pm$0.27 & 0.756 \\
CMJ高度（cm） & 63.28$\pm$6.82 & 62.15$\pm$6.54 & 0.658 \\
SJ高度（cm） & 58.12$\pm$6.35 & 56.89$\pm$5.92 & 0.605 \\
训练自我效能（分） & 80.43$\pm$7.85 & 72.58$\pm$8.12 & 0.012 \\
\bottomrule
\end{tabular}
\begin{tablenotes}[flushleft]\footnotesize
\item 注：ITT样本将所有随机化受试者纳入分析，无论其是否完成干预；脱落受试者按末次观测结转法（BOCF）处理；$p$值来自协方差分析（协变量为基线值与分层变量）。
\end{tablenotes}
\end{threeparttable}
\end{table}
